\section{Parametrization, topology, and trace budgets for the low exponents}
\label{lt:section}\label{sec:low-tools}

The constructions in this Section use length measurements when
$1\leq p<2$, and both length and area measurements when $p=2$.
Area always means parametrized Euclidean area, counted with multiplicity.
The Dirichlet integral uses the squared Frobenius norm, without a factor
$1/2$.  All topological changes are made in the interiors of the open
domains.  In particular, none of the relative statements below concerns
an extension to the boundary of either domain.

The disk and collapse Lemmas convert compatible boundary data and image
measurements into energy bounds for parametrizations.  The topological
Lemmas prepare crossings and exact germs; the simplex squeeze then
converts the counted crossings into trace measurements.  Source and
target sampling control the resulting budgets.

\subsection{Boundary parametrization of a disk}

The conversion of area into nearly minimal Dirichlet energy by a
change of parameters has a classical precursor in Morrey's
$\varepsilon$-conformal parametrization theorem
\cite[Chapter~I, Section~9, Theorem~1.2]{Morrey1948}. Here we also
need exact boundary agreement and control of the boundary energy.

\begin{definition}[Essential tangential compatibility]
\label{lt:compatibility}
Let $D$ be a compact polygonal disk or a convex polyhedral cell and let
$v:D\to\R^d$ be Lipschitz.  We say that $v$ has essential tangential
compatibility if the following properties hold.
\begin{enumerate}
\item At almost every interior point $x$, the a.e. derivative of $v$ has
essential limit $Dv(x)$ as its argument tends to $x$.
\item In radial coordinates based at an interior point, at almost every
boundary parameter the derivatives in boundary directions have essential
limits, from the interior as well, equal to the derivatives of the
boundary restriction.
\end{enumerate}
The exceptional sets here are null sets on the relevant parameter
manifold.  An ambient null-set assertion alone is insufficient for the
second condition.
\end{definition}

\begin{lemma}[Compatibility of finite closed-piece formulas]
\label{lt:closed-pieces}
A continuous Lipschitz map given on a compact parameter manifold by
finitely many $C^1$ formulas extending to neighborhoods of their closed
pieces has the compatibility in Definition~\ref{lt:compatibility}.
The same assertion applies to compositions of such maps.  It continues
to hold after a boundary parametrization is changed on each edge by a
bi-Lipschitz map whose derivative has an essential limit almost
everywhere, and the change is extended by coning.
\end{lemma}

\begin{proof}
At a point belonging to a closed piece, only pieces containing that
point can meet every sufficiently small neighborhood of it.  At almost
every point of an overlap, that point is a density point of the overlap
in its parameter manifold.  Two $C^1$ extensions whose values agree
there have equal derivatives in its tangent directions: their difference
vanishes on a set of density one, and its first-order expansion therefore
vanishes on every tangent vector.  Continuity of the finitely many
extension derivatives now gives the asserted essential limits.
Apply this argument on the boundary parameter manifold for boundary
directions.  For a composition, partition by the finitely many choices
of formulas and use the chain rule; tangential agreement on a preimage
overlap follows from the same density argument.  In a cone extension,
away from its center and face rays the formula is a smooth radial
factor times the boundary formula.  The asserted limits consequently
follow from those of the boundary derivative; the excluded rays and
parameter exceptions have measure zero.
\end{proof}

\begin{lemma}[Boundary-parametrized disk estimate]
\label{lt:disk}
Put $Q=[-1/2,1/2]^2$.  Let $h:Q\to\R^d$ be a bi-Lipschitz embedding
having essential tangential compatibility, and put
\[
 A(h)=\int_Q J_2Dh,\qquad
 L(h)=\int_{\partial Q}|D_t h|^2.
\]
For every $\eps>0$ there is a bi-Lipschitz self-homeomorphism
$\rho:Q\to Q$, equal to the identity on $\partial Q$, such that
\begin{equation}
 \int_Q|D(h\circ\rho)|^2
       \leq C\bigl(A(h)+L(h)\bigr)+\eps,
 \label{lt:disk-weak}
\end{equation}
where $C$ is absolute.

There is also the following sharp form.  Suppose $h_j$ are such
embeddings, $\sup_jL(h_j)<\infty$,
$A(h_j)\leq1+o(1)$, and
$h_j|_{\partial Q}\to\iota|_{\partial Q}$ uniformly, where
$\iota(x)=(x,0)\in\R^d$.  The parametrizations can be chosen so that
\begin{equation}
 \int_Q|D(h_j\circ\rho_j)|^2\leq2+o(1),
 \qquad
 D(h_j\circ\rho_j)\longrightarrow D\iota
       \quad\hbox{in }L^2(Q).
 \label{lt:disk-sharp}
\end{equation}
The bi-Lipschitz constants of the chosen parametrizations may depend on
the individual embedding.  The conclusions transport to polygons
through fixed shape-controlled, bi-Lipschitz piecewise smooth charts;
the sharp square normalization is used when the constant $2$ matters.
\end{lemma}

\begin{proof}
We first majorize the pullback metric by a smooth metric that is
conformally Euclidean on an added collar. Conformal coordinates on a
rectangle then give the weak energy estimate. For the sharp estimate,
coordinate tests force both the rectangle and its boundary correction
to be nearly isometric.

\paragraph{Metric preparation.}
Fix $0<\tau\leq1$, set $S=1+2\tau$ and
$Q^+=SQ$, and use square radial coordinates.  For $\delta>0$, let
$r:Q^+\to Q$ preserve radial directions, carry the inner $Q$ by the
homothety of factor $1-\delta$, and carry the added shell linearly along
each radius onto $Q\setminus(1-\delta)Q$.  Thus $r$ is bi-Lipschitz
for fixed positive $\delta$.  Given $\xi>0$, we may choose $\delta$
and a smooth positive metric $g$ on a neighborhood of $Q^+$ so that
\begin{align}
 g&\geq D(h\circ r)^TD(h\circ r)\quad\hbox{a.e.},
       \label{lt:metric-dominates}\\
 g&=M^2I\quad\hbox{throughout }Q^+\setminus Q,
       \label{lt:metric-shell}\\
 \operatorname{area}_g(Q)&\leq A(h)+\xi,
 \qquad
 \operatorname{area}_g(Q^+\setminus Q)
       \leq C\tau L(h)+\xi,                 \label{lt:metric-area}\\
 \|h\circ r-h\|_{C^0(\partial Q)}&\leq\xi,
 \qquad
 \int_{\partial Q}|D_t(h\circ r)|^2
       \leq CL(h)+\xi.                    \label{lt:metric-trace}
\end{align}
Here a scalar factor in \eqref{lt:metric-shell} is allowed to vary
smoothly in the shell.

We give the details of this preparation because its boundary assertion
is used at the critical exponent.  On the shell the radial derivative
of $h\circ r$ is bounded by $C\delta\Lip(h)/\tau$.  Essential
tangential compatibility implies that the essential suprema of the
perimeter-direction derivatives in shrinking neighborhoods of almost
every boundary parameter tend to the corresponding boundary derivative
norm.  They are bounded by fixed Lipschitz data.  Covering the perimeter
by short overlapping intervals, taking suprema on slightly enlarged
intervals, and using a smooth partition of unity therefore gives scalar
majorants whose squared integrals are at most $CL(h)+\xi$ once
$\delta$ is sufficiently small.  The negligible radial derivative can
be included in these majorants.  Square radial and ordinary angular
coordinates have uniformly bounded distortion in this shell.  Their
conformal metrics consequently have area at most $C\tau L(h)+\xi$.
Choose in addition a typical sufficiently small inner level to obtain
\eqref{lt:metric-trace}.

On the inner square consider the pullback quadratic form in
\eqref{lt:metric-dominates}.  It is bounded, uniformly positive for the
fixed data, and essentially continuous almost everywhere.  On a fine
partition of unity, majorize it by a representative matrix on each
enlarged patch plus its essential oscillation times the identity and a
small positive scalar matrix.  The resulting smooth positive form
dominates the original form.  The oscillations tend to zero almost
everywhere and are bounded, so its area integral tends to that of the
pullback form by dominated convergence.  The latter is the area of
$h$ on $(1-\delta)Q$, hence is at most $A(h)$.  Call this inner majorant $g_0$, and call the scalar shell
majorant $M^2I$.  Choose a finite number $N$ dominating the operator
norms of $g_0$, $M^2I$, and both one-sided pullback forms near
$\partial Q$.  This is possible for the fixed map and the fixed
positive $\delta$.  In a two-sided neighborhood of $\partial Q$ of
thickness $t$ set the metric equal to $NI$.  Blend $g_0$ to $NI$
on the inner side and $NI$ to $M^2I$ on the outer side by smooth
cutoffs supported in a slightly larger neighborhood.  Each blend is
between two forms dominating the relevant pullback form, so it still
dominates that form; the entire outer blend is scalar.  The metric
is constant on an open neighborhood of the interface, so its jets
match across the sides and corners.  Smooth neighborhoods and cutoffs
can be chosen with area $O(t)$, and the added metric area is at most
$CNt$.  Choose $t<\tau/10$ and then $t$ so small that $CNt$ is below
the remaining error allowance.  Derivatives of the cutoff metric
are not charged to any estimate.  Along a sequence, $N$ may diverge,
but $t$ is chosen after $N$ so that $Nt\to0$.  In particular a
fixed-width outer band of the added shell is unaffected.  This proves \eqref{lt:metric-dominates}--\eqref{lt:metric-trace}.

\paragraph{Conformal coordinates.}
By isothermal coordinates and marked-quadrilateral uniformization
\cite[Lemma~7 and Theorem~6]{AhlforsBers1960}\cite[Theorem~4.31]{McMullen2023},
uniformize the smooth metric quadrilateral $(Q^+,g)$ conformally onto
the centered rectangle
\[
 R_{W,S}=[-W/2,W/2]\times[-S/2,S/2],
\]
with corresponding corners, and write this map as $\Phi$.  The marked
quadrilateral uniformization determines the modulus $W/S$; a Euclidean
dilation sets the height to $S$.  In the interior $\Phi$ is smooth and
nonsingular. We justify this regularity before using $\Phi$ in the
energy estimates.
The metric's
Beltrami coefficient $\mu$ is smooth and vanishes on the outer scalar
band, so it extends by zero to a compactly supported smooth coefficient
on the plane. Lemma~7 of Ahlfors--Bers gives a $C^1$ solution with
positive Jacobian. Its smooth bootstrap follows from their
parameter construction \cite[Section~2, Theorem~2]{AhlforsBers1960}:
choose $q>2$ with $\|\mu\|_\infty C_q<1$ for the norm $C_q$ of
their singular integral operator. Translated coefficients
$\mu_t(z)=\mu(z+t)$ vary smoothly in $L^\infty\cap L^q$; the
uniformly invertible operator in that construction therefore gives
smooth dependence of $U_t-\Id$ in its solution space $B_q$, where
$U_t(0)=0$ and $\partial U_t-1\in L^q$. The H\"older norm of $B_q$
makes point evaluation continuous.
Writing $U=U_0$, uniqueness gives $U_t(z)=U(z+t)-U(t)$. On any bounded neighborhood
of $t$, choose a fixed $z_*$ such that $z_*+t$ stays outside
$\supp\mu$. Then
$U(t)=U(z_*+t)-U_t(z_*)$ is smooth, since its first term is
holomorphic there and its second depends smoothly on $t$.
The conformal rectangle map preserves interior smoothness and
nonsingularity. Since the metric is conformally Euclidean near the outer
boundary, Schwarz reflection across sides
\cite[Chapter~4, Section~6.5, Theorem~24]{Ahlfors1979}, and then across the two
perpendicular sides at each corner, shows that $\Phi$ and its inverse
are bi-Lipschitz up to the boundary for each fixed metric.

\paragraph{Weak estimate.}
Fix, for example, $\tau=1/4$. Extend the first
and second Euclidean coordinate functions from the outer boundary
through the shell using cutoffs which vanish before reaching $Q$.
Their $g$-Dirichlet integrals are bounded by an absolute constant:
on the shell the scalar factor cancels from a two-dimensional
Dirichlet integral.  On the rectangle, the corresponding horizontal
and vertical boundary differences equal $S$.  Cauchy's inequality on
horizontal and vertical segments gives, respectively, the lower bounds
\begin{equation}
 S^3/W\quad\hbox{and}\quad SW.
 \label{lt:rectangle-lower}
\end{equation}
It follows that $W$ is bounded above and below by positive absolute
constants.

Choose $\eta<c\tau$, for example $\eta=\tau/16$.
Because $g$ is scalar on the whole added shell, $\Phi$ is ordinary
holomorphic there.  Reflect in each outer source side and its
corresponding rectangle side.  Near a corner use the two perpendicular
reflections, which commute.  Every point in the resulting neighborhood
of width $\eta$ folds into the original shell, so these formulas give
one analytic extension on that entire neighborhood.  Away from its
seams the derivative is an original nonzero derivative or its conjugate.
Across a side the reflected halves map to opposite half-planes, and
at a corner the four copies map into the four distinct target
quadrants.  The extension is therefore locally injective also on
seams and corners, and its derivative is nonzero throughout the band.
Its image is contained in
$[-3W/2,3W/2]\times[-3S/2,3S/2]$, since at most two side reflections
are used.  The rectangle bounds and Cauchy's estimate on disks in the
band now give a uniform upper derivative bound on a slightly narrower
band, independent of the metric on the inner square.

Choose $C$ strictly above that upper bound.  The positive harmonic
function $u=\log(C/|\Phi'|)$ is defined on this fixed-width band.
On each outer side the integral of $|\Phi'|$ equals the corresponding
rectangle side length.  Thus some boundary point has
$|\Phi'|\geq c_0>0$.  A fixed finite chain of disks along the outer
boundary gives by Harnack's inequality
$u\leq C_H\log(C/c_0)$ there.  Consequently
$|\Phi'|\geq C\exp(-C_H\log(C/c_0))>0$ on the entire boundary,
including the corners.  The constants depend on the fixed $\tau$,
but not on the inner metric distortion.  The boundary map of
$\Phi$ thus has a uniformly controlled bi-Lipschitz extension
$P:Q^+\to R_{W,S}$, obtained by coning from the centers.

Set, for $x\in Q$,
\[
 F(x)=h\circ r\circ\Phi^{-1}\circ P(Sx).
\]
Since $P=\Phi$ on $\partial Q^+$ and $r(Sx)=x$ there, this is
$h\circ\rho$ with $\rho$ a bi-Lipschitz self-homeomorphism fixing
$\partial Q$.  The homothety does not change two-dimensional
Dirichlet energy.  Conformal invariance, the bounded distortion of
$P$, and \eqref{lt:metric-dominates} yield
\[
 \int_Q|DF|^2
 \leq C\int_{Q^+}|d(h\circ r)|_g^2\,d\operatorname{area}_g
 \leq 2C\operatorname{area}_g(Q^+).
\]
Letting $\xi$ be sufficiently small proves \eqref{lt:disk-weak}.

\paragraph{Sharp estimate.}
First fix an arbitrarily small positive
$\tau$ and make the metric errors tend to zero along the given
sequence.  Denote the first two physical coordinates of $h_j\circ r_j$
on $Q$ by $X_j,Y_j$.  Each coordinate has squared $g_j$-gradient at
most one, so each has energy at most
$\operatorname{area}_{g_j}(Q)\leq1+o(1)$.  By
\eqref{lt:metric-trace}, its boundary discrepancy from the corresponding
Euclidean coordinate is bounded in $W^{1,2}(\partial Q)$ and tends to
zero in $L^2(\partial Q)$.  The interpolation inequality
\[
 \|v\|_{H^{1/2}(\partial Q)}^2
      \leq C\|v\|_{L^2(\partial Q)}\|v\|_{H^1(\partial Q)}
\]
therefore makes the discrepancy tend to zero in $H^{1/2}$.
The trace extension operator on the fixed polygonal collar, followed
by a cutoff, extends it with vanishing Euclidean $W^{1,2}$ norm and
with support away from the outer half of the shell.  Extend $X_j,Y_j$
accordingly so that they equal the Euclidean coordinates on that outer
half.  Their individual shell energies are at most
$S^2-1+o(1)$, since the metric is scalar there.  Their total individual
energies on $Q^+$ are thus at most $S^2+o(1)$.

Apply \eqref{lt:rectangle-lower} to these tests on
$R_{W_j,S}$.  Both $S^3/W_j$ and $SW_j$ are at most $S^2+o(1)$;
hence $W_j\to S$.  Moreover, if
$(\widehat X_j,\widehat Y_j)=(X_j,Y_j)\circ\Phi_j^{-1}$, equality
asymptotically in the two segmentwise inequalities gives
\begin{equation}
 D(\widehat X_j,\widehat Y_j)\longrightarrow I
              \quad\hbox{in }L^2(R_{W_j,S}).
 \label{lt:coordinate-rigidity}
\end{equation}
For example, the horizontal integral of
$\partial_1\widehat X_j$ on every segment is $S$, and subtracting
$S^3/W_j$ from its energy gives exactly
\[
 \int_{R_{W_j,S}}
 \bigl(|\partial_1\widehat X_j-S/W_j|^2
               +|\partial_2\widehat X_j|^2\bigr).
\]
The vertical coordinate has the analogous identity with derivative
$1$.

On the outer coordinate band, the pair in
\eqref{lt:coordinate-rigidity} equals $\Phi_j^{-1}$.  For an
orientation-preserving similarity $A$ in dimension two,
\[
 |A^{-1}-I|^2\det A=|A-I|^2.
\]
Change of variables therefore gives $D\Phi_j\to I$ in Euclidean
$L^2$ on the corresponding source band.  Reflection and the analytic
interior estimates give uniform derivative convergence on a smaller
band containing the boundary.  The corner normalization and $W_j\to S$
also give convergence of the values there.  Blend $\Phi_j$ in this
fixed band with the linear map
$(x_1,x_2)\mapsto(W_jx_1/S,x_2)$.  This gives extensions $P_j$ of the
boundary values with distortion $1+o(1)$; for large $j$ their derivatives
are close to the identity, so they are bi-Lipschitz homeomorphisms onto
the rectangle.  The preceding energy argument now gives
\[
 \int_Q|D(h_j\circ\rho_j)|^2
   \leq(1+o(1))\,2\operatorname{area}_{g_j}(Q^+)
   \leq2+C\tau\sup_jL(h_j)+o(1).
\]
Take $\tau\downarrow0$ diagonally.  Finally, boundary agreement and
uniform convergence of the boundary values give, by integration in
coordinate directions,
\[
 \int_Q D(h_j\circ\rho_j):D\iota\longrightarrow2.
\]
Expanding the square of the derivative difference and using the energy
bound proves the second conclusion in \eqref{lt:disk-sharp}.
\end{proof}

\subsection{Subcritical collapse onto boundary data}

The skeleton and radial-extension strategy is classical in Sobolev
approximation; see Bethuel \cite[Section~II.1, pp.~164--165, and
Section~II.3, Eqs.~(32), (37), pp.~170--171]{Bethuel1991}.
This background motivates the collapse below.  Its injective
reparametrization, fixed image and exact boundary compatibility are
proved locally; Bethuel's density theorem does not supply these
homeomorphism requirements.

\begin{lemma}[Cell collapse]
\label{lt:collapse}
Let $m\in\{2,3\}$ and $1\leq p<m$.  Let $D\subset\R^m$ be a
convex polyhedral cell of scale $\ell$, with inradius bounded below by
$c\ell$ and diameter bounded above by $C\ell$.  Suppose
$H:D\to\R^d$ is a bi-Lipschitz embedding with essential tangential
compatibility.  Let $\sigma:\partial D\to\partial D$ be a
face-preserving, cellwise bi-Lipschitz homeomorphism.  For every
$\eps>0$ there is a bi-Lipschitz parametrization $F:D\to H(D)$ with
$F|_{\partial D}=H\circ\sigma$ and
\begin{equation}
 \int_D|DF|^p
       \leq C_{m,p,c,C}\ell
                    \int_{\partial D}|D_t(H\circ\sigma)|^p+\eps.
 \label{lt:collapse-estimate}
\end{equation}
The boundary data on shared cells are retained exactly, so the
parametrizations glue when those data agree.
\end{lemma}

\begin{proof}
Rescale to $\ell=1$ and place an incenter at the origin. Cone-extend
$\sigma$ from that center; this is bi-Lipschitz for the fixed data.
In radial coordinates
$x=rz$, $z\in\partial D$ and $0\leq r\leq1$, precompose the
resulting map by the radial homeomorphism with profile
\[
 q_{\alpha,\delta}(r)=
 \begin{cases}
 (1-\delta)r/\alpha,&0\leq r\leq\alpha,\\
 1-\delta+\delta(r-\alpha)/(1-\alpha),&\alpha\leq r\leq1.
 \end{cases}
\]
Here $0<\alpha,\delta<1$, and the resulting parametrization is
\[
 F_{\alpha,\delta}(rz)
       =H\bigl(q_{\alpha,\delta}(r)\sigma(z)\bigr).
\]
On $r\leq\alpha$ all derivatives are
at most a fixed-data constant times $\alpha^{-1}$; hence the energy
there is $O(\alpha^{m-p})$.  On $r>\alpha$ the radial derivative
tends to zero as $\delta\downarrow0$, while the tangential derivatives
tend essentially to $r^{-1}D_t(H\circ\sigma)(z)$. Indeed, the
angular derivative contains
$q_{\alpha,\delta}(r)DH(q_{\alpha,\delta}(r)\sigma(z))D_t\sigma(z)$.
The last factor is evaluated at the fixed boundary point $z$;
only the base map's tangential derivatives require a limiting
argument. The radial
Jacobian is comparable to $r^{m-1}$ with constants determined by the
cell shape.  Essential compatibility and dominated convergence yield
an upper bound
\[
 C\left(\int_\alpha^1r^{m-1-p}\,dr\right)
                      \int_{\partial D}|D_t(H\circ\sigma)|^p
\]
for the limiting outer energy.  To apply the essential limits, use a
sequence of the individual bi-Lipschitz radial changes and the common
full-measure set of their Lebesgue representatives.  Coning $\sigma$
only inserts its fixed a.e. angular derivative.  Since $p<m$, the
radial integral stays bounded as $\alpha\downarrow0$ and the central
error tends to zero.  First choose $\alpha$ and then $\delta$.
Restoring scale multiplies the boundary integral by $\ell$ and proves
\eqref{lt:collapse-estimate}.  All radial profiles fix the boundary,
which proves the gluing assertion.
\end{proof}

\subsection{Local topological modifications}\label{lt:local-topology}

The budget construction needs regular transverse crossings at its
measured hits, while the radial blocks need exact identity germs at
selected centers. Both are local topological requirements; neither
provides an energy estimate. We first treat one crossing. We then
control the marked rays meeting at a center, where the way they join
across a surrounding shell must also be retained.

All constructions in this subsection take place in interior coordinate
neighborhoods of three-manifolds. The classical inputs are the locally
flat Schoenflies theorem, Dehn's lemma, and the classification of
mapping classes of punctured spheres; see
\cite{Brown1960,Brown1962,Papakyriakopoulos1957,FarbMargalit2012}.
We use the relative approximation and small embedding extension
statements of Lemmas~\ref{sm:relative-pl} and
\ref{sm:local-extension}. Marked points on a sphere carry no tangent
framing; an isotopy fixing such a point need not fix a tangent
direction there.

\begin{lemma}[Cleaning an isolated crossing]\label{lt:clean-crossing}
Let $P$ be a locally flat embedded surface disk and let $I$ be an interval
which is straight in specified local topological coordinates.  Suppose
that $z\in P\cap I$ is interior to both objects, is an isolated point of
their intersection, and has a neighborhood containing no other surface
sheet under consideration.  There is an ambient modification of $P$,
supported in an arbitrarily small neighborhood of $z$, after which the
intersection in that neighborhood is either empty or consists of a
single standard transverse crossing.  At a retained crossing the disk
can be taken flat across the straight interval.  No initial local
flatness of the pair $(P,I)$ is assumed.

If the interval passes from one local side of $P$ to the other at $z$,
then a sufficiently small modification retains one crossing.
\end{lemma}

\begin{proof}
Choose a small topological ball $B$ about $z$ in which $P$ is a proper
unknotted disk.  All changes will occur away from $\partial B$.
Straighten $I$ in its supplied coordinates near $z$ and choose two
cap levels on opposite sides of $z$ within a segment of the axis whose
only intersection with $P$ is $z$.  In a parameter disk for $P$ choose
round disks $D_1\Subset\operatorname{int}D_0$ about the parameter of
$z$ so that $P(D_0)$ is compactly contained in the open slab between
these levels and in the straight-axis chart.  Now choose a sufficiently
thin polygonal cylinder $C$ about the intervening axis interval.  Its
caps miss $P$.  Compactness and the isolated-axis intersection make
all intersections with its side annulus $A$ lie in $P(D_1)$ once
the radius is sufficiently small.  Every side loop and its bounded
parameter subdisk then lie in $D_1$.  Thus all disks subsequently used
for linking lie inside the cap slab and cannot meet the axis extended
beyond the isolated supplied segment.

Put $P$ in PL general position near $A$, without changing it near the
axis or near $\partial B$.  To justify this preliminary operation,
write $P$ as the image of a flat disk under a chart homeomorphism,
approximate that homeomorphism by a PL homeomorphism on a compact
neighborhood of the prospective side intersections, and install the
approximation on a smaller closed neighborhood by
Lemma~\ref{sm:local-extension}.  The support misses the axis, and the
unsupported part of the disk has a positive gap from the cutting side.
Take the perturbation sufficiently small that $P(D_0)$ remains in
the open cap slab and all prospective side intersections remain in
its parameter interior.  A small change of the polygonal side gives
transverse intersections.
Consequently $P\cap A$ is a finite family of polygonal simple closed
curves.

First remove every curve which is null-homotopic in $A$.  Choose an
innermost such curve on $A$.  Its annular disk has interior disjoint
from $P$: it contains no further null curve by the choice, and cannot
contain an essential curve of $A$.  Replace the disk which this curve
bounds in $P$ by that disk on $A$, pushed slightly off $A$.  Round the
seam on the side dictated by the remaining part of $P$.  The annular
disk lies inside the open cap slab; its small push can be kept there.
The operation replaces a parameter subdisk of $D_0$, and hence the
remaining relevant subdisks still lie in that slab.  It gives an
embedded locally flat disk, removes the chosen intersection curve, and
creates no new side or axis intersections.  Repeating finitely many
times leaves only essential side curves.  The disks discarded in this
operation may contain $z$; in that event the axis intersection is simply
removed.

Every remaining essential side curve has linking number $\pm1$ with
the extended straight axis.  Its disk in $P$ must therefore meet that
axis.  Since the surgeries have added no axis intersections, this disk
contains the original isolated intersection.  The disks bounded by the
remaining curves are thus nested.  Replace the outermost one by a tame
proper disk inside $C$ with the same rim and with exactly one flat
transverse crossing of the axis.  Such a disk is obtained by isotoping
the essential rim in the side annulus to a circular section and
extending that isotopy across an outer collar of a meridian disk.
After this replacement the rest of $P$ is outside the cylinder.  If no
side curves remain, there is no intersection with the cylinder at all:
the caps miss $P$, the outer boundary of the proper disk lies outside
$C$, and any component entering $C$ would have to cross its side.

The result is a locally flat proper disk in $B$ agreeing with the
original disk near its boundary.  Both disks split $B$ into balls.
The locally flat Schoenflies theorem \cite[Theorem~4]{Brown1962},
followed by extension on the two
sides, gives an ambient homeomorphism of $B$, fixed on $\partial B$,
carrying the original disk to the constructed disk.  Extend it by the
identity outside $B$.  The ball, the cylinder, and each surgery support
can be chosen successively smaller, proving the support assertion.

Finally, in the opposite-side case retain two nearby axis points on
opposite sides, outside the modification support.  Separation forces
the modified disk to meet the intervening interval.  Since the
construction leaves at most one intersection, exactly one remains.
\end{proof}

\begin{corollary}[Affine germs at cleaned crossings]
\label{lt:affine-crossing}
Let $h$ be a homeomorphism between open three-manifolds.  A crossing of
a regular source surface with the preimage of a target line can be
cleaned by applying Lemma~\ref{lt:clean-crossing} to the image surface.
A crossing of a regular source curve with the preimage of a target
surface can be cleaned by applying the Lemma to $h^{-1}$.
If the regular strata are straightened in $C^1$ coordinates, the
modified homeomorphism can additionally be made affine on a smaller
neighborhood of every retained crossing.  The only crossings there
are the retained transverse ones.  Supports can be prescribed to be
small in both the source and target.
\end{corollary}

\begin{proof}
After cleaning, use the indicated coordinates so that the relevant
source and image surface germs are planes.  Choose a small outer source
ball in this germ and a much smaller inner source ball.  Prescribe an
affine map on the inner ball, sending its equator to the target plane,
with its image contained in the outer image ball.  Choose the affine
map to agree with the side correspondence and orientation of $h$.

Extend first across the planar annulus between the two equators.
On either side, the outer half-ball is a topological ball with a flat
boundary patch near the equatorial center.  Removing the inner
half-ball, attached along a disk in that patch, again leaves a ball.
This follows either from a boundary half-space chart and relative
Schoenflies or from the usual ball-with-boundary-attached-ball model.
The resulting boundary homeomorphisms extend across these balls.
The two extensions agree on the planar annulus, so they define the
required cutoff of $h$.  The surface remains in the target plane during
this further change.  In the inverse construction the inverse affine
germ is again affine.  Finally, continuity of $h$ and $h^{-1}$ permits
the paired support neighborhoods to be made small in both spaces.
\end{proof}

A crossing has now been made regular without changing the map away
from a small paired neighborhood. To install an identity germ at a
center with several prescribed rays, we must also control the arcs
between successive surrounding spheres and their boundary
identifications. The next lemma identifies the individual arcs; the
following product lemma identifies the shell together with all its
labelled strands. These are the two topological inputs to
Proposition~\ref{lt:star-replacement}.

\begin{lemma}[A two-point cut of an unknot]\label{lt:unknot-cut}
Let $K$ be a tame unknot in the three-sphere and let $B$ be a locally
flat ball whose boundary meets $K$ in exactly two standard transverse
crossings.  Then the proper arc $K\cap B$ is boundary-parallel in $B$.
In particular, the ball--arc pair is the standard unknotted interval
pair.
\end{lemma}

\begin{proof}
Choose a tube about $K$ meeting $\partial B$ in two meridian disks.
Such tubes can be constructed on the two cut intervals, starting with
the product charts at their endpoints and continuing ordinary regular
neighborhoods along their interiors.  Glue the endpoint disks to obtain
the tube about $K$.  The exterior of the whole tube has fundamental
group $\mathbb Z$, generated by a meridian.  Its two pieces are glued
along the annulus obtained from $\partial B$ by deleting the meridian
disks.  The inclusion of this annulus into the whole exterior induces
an isomorphism on fundamental groups.  Hence its inclusions into the
two pieces are injective.  The injective-amalgam form of van Kampen
then embeds each piece's group in the total group.  Since its image
contains the meridian generator, each piece has group $\mathbb Z$,
again generated by that meridian.

In the exterior piece lying in $B$, join the two end circles by an arc
on the annulus from $\partial B$, and close it by a connector on the
lateral annulus of the drilled tube.  This is a simple loop on the
boundary of the exterior.  Its class is a power of the meridian.
Twisting the connector by the opposite integer power on its annulus
makes the loop null-homotopic while preserving simplicity.  Dehn's
Lemma supplies a properly embedded disk with that boundary.

Restore the tube and attach the radial strip joining its connector to
the core arc.  The strip is embedded even when the connector twists:
in product coordinates its angle may vary with the interval coordinate
while its radius decreases to zero at the core.  Its two short edges
lie in the meridian end disks.  The resulting locally flat disk has
boundary consisting of $K\cap B$ and an arc in $\partial B$.  Thus the
arc is boundary-parallel.  Sliding along a neighborhood of this disk
identifies the pair with a ball and a standard diameter.  Rounding the
seams does not change this conclusion.
\end{proof}

\begin{lemma}[A sphere crossing radial strands once]\label{lt:punctured-product}\label{lt:ray-product}
Let $0<a<b$, let $v_1,\ldots,v_n\in S^2$ be distinct, with $n\geq3$, and put
\[
 M=\{x\in\R^3:a\leq |x|\leq b\},\qquad
 L_i=\{t v_i:a\leq t\leq b\}.
\]
Suppose that $S\subset\operatorname{int}M$ is a locally flat embedded
sphere separating the two boundary spheres of $M$.  Suppose that $S$
meets each $L_i$ at exactly one point $z_i$, and that these intersections
are standard topological transverse crossings: the surface--arc pair
has a local chart onto a plane and a transverse straight line.
Let $U$ be the closed region between $\{|x|=a\}$ and $S$.
Then there is a homeomorphism of pairs
\[
 \Phi:\bigl(S^2\times[0,1],\{v_1,\ldots,v_n\}\times[0,1]\bigr)
 \longrightarrow
 \bigl(U,\textstyle\bigcup_i(L_i\cap U)\bigr)
\]
with $\Phi(q,0)=a q$ and $\Phi(v_i,1)=z_i$.
Moreover, for any such product parametrization, the upper boundary map
$\phi(q)=\Phi(q,1)$ satisfies
\[
 q\longmapsto \frac{\phi(q)}{|\phi(q)|}\simeq\Id_{S^2}
 \quad\text{relative to }\{v_1,\ldots,v_n\}.
\]
This homotopy takes unmarked points to unmarked points at every time.
Equivalently, the link identification supplied by the product agrees,
up to a homotopy preserving the marked points and their complement,
with radial projection.
\end{lemma}

\begin{proof}
Write $P=S^2\setminus\{v_1,\ldots,v_n\}$ and
$S^\circ=S\setminus\{z_1,\ldots,z_n\}$.  Radial projection restricts to
\[
 f:S^\circ\longrightarrow P.
\]
This map is proper: its extension to the compact sphere $S$ has
$f^{-1}(v_i)=\{z_i\}$, so the inverse image of a compact subset of $P$
is a compact subset of $S^\circ$.  Orient $S$ as the boundary of its
bounded complementary region.  Since $S$ encloses the inner sphere,
its unpunctured radial projection has degree $1$.  Local degree at any
point of $P$ therefore shows that the proper map $f$ also has degree $1$.

We record explicitly the covering argument for $f_*$.  Let
$p:\widetilde P\to P$ be the connected covering corresponding to
$f_*\pi_1(S^\circ)$, and let $\widetilde f$ be the lift of $f$.
The lift is proper: for compact $K\subset\widetilde P$, its inverse
image is closed in the compact set $f^{-1}(p(K))$.
Give the covering its induced orientation and let $d$ be the proper
degree of $\widetilde f$.  For a fixed $y\in P$, compactness of
$f^{-1}(y)$ implies that its lifted image meets only finitely many
points of the discrete closed fibre $p^{-1}(y)$.  Additivity of local
degree expresses $\deg f$ as the sum of the local degrees over these
points.  At every point of the covering the local-degree sum for
$\widetilde f$ equals $d$; it is zero at a point outside its image.
Consequently an infinite-sheeted covering would force $d=0$ and then
$\deg f=0$.  For a finite-sheeted covering with $k$ sheets, the same
calculation gives $1=k d$.  Hence $k=1$, and $f_*$ is surjective.
Both fundamental groups are free of rank $n-1$, so the Hopf property
of finitely generated free groups \cite[III.A, Complements~18--19]{deLaHarpe2000} makes $f_*$ an isomorphism.

Choose pairwise disjoint closed regular-neighborhood tubes $N_i$
around the full strands $L_i$, each meeting $S$ in one meridian disk
and meeting each boundary sphere of $M$ in one end disk.  Such tubes
are obtained from the pair charts at the crossings and regular
neighborhoods of the two remaining tame subintervals, using the same
end disks where they join.  In particular all tubes and their pieces
on either side of $S$ are standard interval neighborhoods.
Drill out their interiors relative to $M$, obtaining an exterior $E$.
Denote its bottom and top planar boundary surfaces by $F_0$ and $F_2$,
and the truncated copy of $S$ by $F_1$.
Regular-neighborhood coordinates identify the fundamental groups of
these exteriors with those obtained by deleting just the core
strands, compatibly with the inclusions of the truncated surfaces.
Since
\[
 M\setminus\bigcup_iL_i\cong P\times[a,b],
\]
our calculation of $f_*$ shows that $F_1\hookrightarrow E$ induces a
fundamental-group isomorphism.  The same is immediate for $F_0$ and
$F_2$.

Cut $E$ along $F_1$, and let $W$ be the piece on the inner side.
The inclusion of $F_1$ into either piece is injective on fundamental
groups, because its composite into $E$ is injective.  Van Kampen and
the normal-form Theorem for a free product with amalgamation therefore
embed both piece groups in $\pi_1(E)$.  As the common group
$\pi_1(F_1)$ already maps onto $\pi_1(E)$, both piece inclusions are
isomorphisms.  It follows that both
\[
 \pi_1(F_0)\longrightarrow\pi_1(W),\qquad
 \pi_1(F_1)\longrightarrow\pi_1(W)
\]
are isomorphisms.  Apart from these horizontal surfaces, the boundary
of $W$ consists of $n$ annuli $A_i$, one from each drilled tube, joining
the corresponding boundary circles of $F_0$ and $F_1$.

The manifold $W$ is irreducible.  Indeed, let $\Sigma$ be a locally
flat sphere in its interior.  It cannot separate the boundary spheres
of $M$, since every full strand joins those spheres and is disjoint
from $\Sigma$.  Schoenflies therefore gives a ball $B$ bounded by
$\Sigma$ and contained in $\operatorname{int}M$.  Each $N_i$ is connected,
is disjoint from $\Sigma$, and meets $\partial M$, so it lies outside
$B$.  The connected surface $F_1$ also lies outside $B$, since it is
disjoint from $\Sigma$ and its boundary meets the tubes.  Thus
$B\subset W$.  All subsequent disks and spheres may be taken locally
flat, or PL after triangulating this compact $3$-manifold.

We now give a relative product recognition argument, including the
boundary identifications.  Orient $F_0$ and $F_1$ so that oriented
peripheral loops correspond across the annuli $A_i$; these are opposite
to one another when compared with their induced orientations as parts
of $\partial W$.  The fundamental-group identification between them
preserves the individual oriented peripheral conjugacy classes.
The surface classification Theorem in its peripheral-preserving
Dehn--Nielsen form supplies a homeomorphism
$j:F_0\to F_1$ realizing this outer isomorphism and matching the labelled
boundary circles \cite[Theorem~8.8 and Proposition~3.19]{FarbMargalit2012}.

Here is the adjustment from outer to based identifications.  For this
step relabel the boundary circles and their annuli by $0,\ldots,n-1$.
Choose
a basepoint $x_0$ on one boundary circle $C_0$ of $F_0$, and a connector
$\Delta_0$ in its annulus from $x_0$ to $j(x_0)$.  Transport the second
inclusion into $W$ along $\Delta_0$.  Its composition with $j_*$ and
the first inclusion differ by an inner automorphism.  Both carry the
positive based peripheral of $C_0$ to the same element $g_0$ of
$\pi_1(W,x_0)$, so the conjugating element centralizes $g_0$.
Every peripheral of an $n$-holed sphere is primitive in its free
fundamental group; since $n\geq3$, its centralizer is precisely the
cyclic group it generates.  Twisting $\Delta_0$ by an appropriate
integer power in its annulus therefore gives the based equality
\[
 [c]=[\Delta_0\,j(c)\,\Delta_0^{-1}]
 \quad\text{in }\pi_1(W,x_0)
 \quad\text{for every based loop }c\text{ in }F_0.
 \tag{*}
\]

Choose disjoint properly embedded arcs $\alpha_1,\ldots,\alpha_{n-1}$
in $F_0$, joining $C_0$ to the other boundary circles, whose cutting
opens $F_0$ to a disk.  Their endpoints on $C_0$ are distinct.
For each other circle choose the endpoint $x_i$ of its arc, a path
$c_i$ from $x_0$ to $x_i$ following $C_0$ and then $\alpha_i$, and an
initial connector $\Delta_i$ in $A_i$ from $x_i$ to $j(x_i)$.
The two paths
\[
 c_i\Delta_i\quad\text{and}\quad\Delta_0j(c_i)
\]
transport the peripheral at $j(x_i)$ to the same based element of
$\pi_1(W,x_0)$, by (*) and the annulus correspondence.
Their discrepancy thus centralizes that peripheral.  The same
centralizer argument allows an integer twist of $\Delta_i$ making
these two paths homotopic relative to their endpoints in $W$.

Parametrize each $A_i$ as a circle product, with its lower boundary
the identity, its upper boundary the restriction of $j$, and its
chosen connector in the adjusted relative homotopy class.  Such a
parametrization exists because its remaining integer winding is
changed by a Dehn twist of the annulus.  On the base annulus use the
parametrization with track $\Delta_0$ at $x_0$ and use its tracks also
at every endpoint of an $\alpha_i$ on $C_0$.  Thus all connectors are
disjoint.  The product tracks are compatible with travel along a
boundary circle, so the preceding path equalities imply that each
closed curve formed from $\alpha_i$, its two connector tracks, and
$j(\alpha_i)$ is nullhomotopic in $W$.  These curves, denoted $\gamma_i$,
are disjoint simple curves on $\partial W$.

Dehn's Lemma supplies properly embedded disks $D_i\subset W$ with
$\partial D_i=\gamma_i$ \cite{Papakyriakopoulos1957}.  They can be chosen
disjoint: put them in general position and remove their intersection
circles by the usual innermost-disk replacements, keeping their
already disjoint boundaries fixed.  The chosen horizontal and annular
maps give a boundary homeomorphism from $\partial(F_0\times[0,1])$
to $\partial W$ carrying the boundaries of the rectangles
$\alpha_i\times[0,1]$ to the $\gamma_i$.  Consequently cutting $W$
along the $D_i$ produces a manifold whose boundary is one sphere,
just as cutting $F_0\times[0,1]$ along those rectangles does.
Every component of a manifold obtained by these proper disk cuts has
nonempty boundary, so the cut manifold is connected.  Irreducibility
persists under the cuts: a sphere in the cut manifold bounds a ball
in $W$, and each cutting disk, being connected and having boundary
on $\partial W$, must lie outside that ball.  A collar inset of the
spherical boundary then shows that the cut manifold is a ball.

Extend the boundary correspondence over the rectangles and the
$D_i$, and then over the resulting balls.  Regluing gives a product
homeomorphism $F_0\times[0,1]\to W$ realizing the chosen boundary
correspondences.  Restore each removed tube piece.  Extend its
boundary-circle maps over its two cap disks, taking the marked core
endpoints to one another.  A homeomorphism of the boundary of a
standard ball--interval pair preserving the two endpoints extends to
the pair: identify both pairs with a ball and a diameter and cone the
boundary map.  Applying this to each tube extends the exterior
product to the asserted product of the filled region and its labelled
strands.  At the lower end choose the cap extensions to agree with
the identity, so that $\Phi(q,0)=a q$.

Finally the formula
\[
 (q,t)\longmapsto\frac{\Phi(q,t)}{|\Phi(q,t)|}
\]
is the claimed boundary-comparison homotopy.  It fixes each $v_i$
because its product track is contained in $L_i$, and it avoids all
marked directions whenever $q$ is unmarked because the product
homeomorphism is a homeomorphism of pairs.
\end{proof}

\begin{proposition}[Replacement at a star center]\label{lt:star-replacement}\label{lt:star}
Use centered Euclidean coordinates in the source and target of a local
homeomorphism $h$, with $h(0)=0$.  Let
\[
 D=\{v_1,\ldots,v_n\}\subset S^2,\qquad n\geq3,
 \qquad R_i=\{t v_i:t\geq0\},
\]
be distinct labelled directions, partitioned into forward and reverse
sets $F$ and $E$.  Each nonempty set has at least two elements.
Choose narrow round cones $C_i$ about $R_i$, and slightly wider round
cones $C_i^+$ with pairwise disjoint angular closures.  Assume, as germs
at the origin, that
\[
 h(R_i)\subset\operatorname{int}C_i\quad(i\in F),
 \qquad h^{-1}(R_i)\subset\operatorname{int}C_i\quad(i\in E).
\]
For each nonempty distorted system, assume there are nested round
spheres with radii tending to zero in its own space, each meeting every
arc of that system exactly once, with the arc passing from one side of
the sphere to the other.  Impose one of the following angular tests.
\begin{enumerate}
\item There is a finite based generating loop system
$(\gamma_1,\ldots,\gamma_{n-1})$ for $S^2\setminus D$, chosen outside
the closed angular caps of all the cones.  At arbitrarily small common
source radii $r$, all parametrized loops
\[
 s\longmapsto \frac{h(r\gamma_a(s))}{|h(r\gamma_a(s))|}
\]
are close to $\gamma_a$ within a fixed simultaneous homotopy tolerance
in $S^2\setminus D$.  The tolerance also keeps the image loops outside
the cone caps.  Their common basepoint is identified with the original
basepoint through the small basepoint neighborhood in this tolerance.
\item There are only forward directions, all lying cyclically on a
source equator and corresponding to the same target equator
directions.  The map $h$ preserves orientation.  At arbitrarily small
source radii the parametrized radial projection of the image of that
equatorial circle is uniformly close to its original parametrization,
with tolerance small compared with the separations of successive
marked points.
\end{enumerate}
The cones are taken narrow enough relative to these tolerances and
separations.  Then $h$ can be changed in arbitrarily small paired
neighborhoods of the center to equal $\Id$ as a germ.  The forward and
reverse angular conditions persist, with the margins from $C_i$ to
$C_i^+$ if necessary.  Any additional strict safety conditions away
from the germ are preserved by making the supports sufficiently small.
The assertion imposes no identity condition at other points of the
original intervals.
\end{proposition}

\begin{proof}
We first straighten the distorted arc systems within their separate
angular sectors.  Each distorted arc is tame as an ambient arc, being
the image or preimage of a straight interval under a homeomorphism.
At each of its intersections with a chosen round section, straighten
the arc in an ambient chart and apply Lemma~\ref{lt:clean-crossing} to
the section disk.  The opposite-side crossing persists once if the
support is sufficiently small.  Move the changed section back to the
round section by the inverse supported homeomorphism, dragging the
arc with it.  The resulting arc has a standard transverse crossing
with the original round section.  Choose these changes in disjoint
annular patches inside the relevant angular sectors.  They can be
made arbitrarily small and preserve nesting of the sections.  All
section crossings are now standard pair crossings.

Consider the connecting strand between two successive sections inside
one angular sector.  Their annular-sector region is a ball, and this
strand is a proper interval in it.  It is unknotted for the following
reason.  Choose a second arc of the same distorted system, which is
possible by the cardinality hypothesis.  On the undistorted side of
the system, the two radial arcs can be closed outside a smaller germ
neighborhood by a polygonal path to give an unknot.  Take the closing
polygon and its spanning disk in a sufficiently small original chart
ball whose image remains in the working chart.  The image of the
closed curve is still an ambient unknot: the chart-ball homeomorphism
extends across its exterior by locally flat Schoenflies.  For a
reverse system use $h^{-1}$ in this argument.  Choose the round
sections so close to the center that they avoid the closing path.
The second arc lies in its own disjoint sector, so the first sector's
annular ball meets this unknot in just the strand under consideration,
with exactly two boundary crossings.  The section-cleaning
homeomorphisms preserve the ambient unknot.  Lemma~\ref{lt:unknot-cut}
therefore makes the strand boundary-parallel in its annular-sector
ball.

Straighten each such strand in that ball, using matching maps on the
cap disks and the identity on the sector sides.  The prescribed cap
maps present no extra obstruction: for the standard ball--diameter
pair, a homeomorphism of the boundary sphere respecting the two ends
extends by coning, and fixes the diameter as a set.  At the outermost
section make a transition inside the same sector and extend by the
identity outside a slightly larger neighborhood.  The infinitely
many annular straightenings glue to a homeomorphism at the center,
with continuous inverse there, because their supports lie in balls
with radii tending to zero.

Perform the target changes for the forward system and the source
changes for the reverse system.  All the sectors are disjoint, so
the target changes leave the target reverse rays fixed, and the
source changes leave the source forward rays fixed.  Pre- and
post-composing accordingly gives a homeomorphism $H$ carrying every
labelled ray of the combined system to itself as a set, as a germ.
The generator test is unchanged: its source loops and their image
directions have positive margins from the supports.  In the
equatorial case the target changes have arbitrarily small angular
displacement, since they act in sufficiently narrow sectors; thus
the parametric angular test persists within its homotopy tolerance.

Choose a test radius $r$ inside this exact-ray germ.  The locally flat
sphere $S=H(\partial B_r)$ meets each marked ray once in a standard
topological crossing.  Indeed $H$ carries the standard source
sphere--ray pair to this pair.  A sufficiently small round target
sphere lies inside $H(B_r)$; together with $S$ it bounds a shell to
which Lemma~\ref{lt:punctured-product} applies.  Its product comparison
of the marked boundary spheres is homotopic to radial projection on
their punctured complements.

In the first angular case the comparison of the outside boundary map
with the intended identity map induces the identity on the chosen
based generators.  The punctured-sphere mapping-class Theorem implies
that it is isotopic to the identity through maps fixing the labelled
marked points \cite[Theorem~8.8]{FarbMargalit2012}.
The same Theorem detects the orientation here; the
assumption $n\geq3$ excludes the two-puncture ambiguity.

In the equatorial case take a consecutive spanning chain of equatorial
joining arcs between marked points, omitting one closing arc. Their
projected images represent the same
classes relative to their marked endpoints.  To see this directly,
the middle of each parametrized image lies in a narrow corridor
containing no puncture, while each short end lies in a disk containing
only its own marked point.  In such an end disk, possible angular
winding can be interpolated using lifted polar angles on the open
end of the arc; the radius tends to zero at the endpoint throughout
the interpolation.  Thus no framing obstruction remains at that
marked point.  Interior points of the image arcs avoid all marked
points because $H$ already maps the complete labelled ray germs
exactly to themselves.  The homotopy comparison in
Lemma~\ref{lt:punctured-product} transfers these arc classes to the
boundary identification.  The usual disjoint-arc isotopy test fixes
these classes simultaneously
\cite[Section~1.2.7, Proposition~2.8, and Lemma~2.9]{FarbMargalit2012}.
Cutting along this chain leaves a disk,
on which the orientation-preserving boundary identification is
isotopic to the identity.  This proves the required trivial marked
mapping class also in the second case.

Choose a smaller source sphere $\partial B_s$ whose intended identity
image is contained in $H(B_r)$.  Product coordinates on the source
annulus and on the target annulus, together with the preceding
marked-sphere isotopy, extend the outside values of $H$ and the
inside identity across the annulus while respecting every labelled
ray.  Set the map equal to the identity on $B_s$ and equal to $H$
outside $B_r$.  This is a homeomorphism and supplies the asserted
identity germ.  The earlier sector straightenings and the final
cutoff have arbitrarily small support in the corresponding spaces.
Their angular bounds and the given strict margins prove the
remaining assertions.
\end{proof}

\begin{lemma}[Relative regularization with protected germs]
\label{lt:relative-germs}
Suppose a homeomorphism has been made affine in specified $C^1$
coordinates near a disjoint locally finite family of crossing
points, and has prescribed exact affine or identity germs at
finitely many star centers.  It admits a fine approximation by a
locally bi-Lipschitz homeomorphism retaining smaller exact germs.
On each compact the approximation is given by finitely many $C^1$
formulas on closed pieces.  Positive-gap avoidance conditions,
the retained transverse crossing counts, and strict angular
conditions can be preserved.  Fine approximation tolerances can
also impose reciprocal closeness on compact sets.
\end{lemma}

\begin{proof}
First realize the required coordinate changes by supported $C^1$
diffeomorphisms.  After separating an affine first jet, such a local
coordinate change has the form $g(0)=0$, $Dg(0)=I$.  On a sufficiently
small ball, a cutoff
\[
 \widetilde g(x)=x+\chi(x)\bigl(g(x)-x\bigr)
\]
equals $g$ on a smaller ball and the identity outside the larger one.
Because $Dg-I=o(1)$ and $g(x)-x=o(|x|)$ at the center, its derivative
is uniformly close to $I$ when the balls are sufficiently small.
It is therefore a supported $C^1$ diffeomorphism.  Use disjoint
locally finite support neighborhoods in the source and target;
the affine parts are absorbed into the prescribed affine germs.
After these coordinate changes the homeomorphism is literally
affine on neighborhoods of all protected points.

Enclose smaller germs by closed polyhedra inside those affine
neighborhoods and apply the relative PL approximation of
Lemma~\ref{sm:relative-pl}.  Equivalently, excise these neighborhoods
and approximate on their open complement as a PL three-manifold
with boundary, relative to its already prescribed boundary collar.
Only finitely many protected pieces occur on any compact, so the
construction is locally finite.  Undoing the supported $C^1$
coordinate changes gives a locally bi-Lipschitz homeomorphism
with the stated finite piecewise-$C^1$ description on compacts.

The exact germs preserve the crossings there.  Elsewhere the
specified avoidances have positive gaps on the compact pieces
where they are needed, so sufficiently fine approximation creates
no additional intersections.  The same choice preserves angular
inequalities with a strict margin.  Finally, for a homeomorphism
fine source-dependent approximation tolerances can be chosen to
ensure prescribed inverse tolerances on target compacts; use
continuity of the original inverse and compact exhaustion.
\end{proof}


We can now regularize finitely many crossings on each compact set
and retain exact germs at the protected centers. The resulting map
is locally bi-Lipschitz, but its derivative bounds may be arbitrarily
large. The next construction therefore estimates selected image
curves and surfaces: generic sampling controls their intersection
counts, and a target reparametrization converts those counts into
length and area.

\subsection{Generic traces and target triangulations}

\begin{lemma}[Generic trace properties]
\label{lt:trace-slicing}
Let $f:\Omega\to\Omega'$ be a homeomorphism of open subsets of
$\R^3$ with $f\in W^{1,p}_{\mathrm{loc}}$, $1\leq p\leq2$.
Consider locally finite families of compact Lipschitz semialgebraic
parametrizations of dimension $m=1$, and also of dimension $m=2$ if
$p=2$.  Almost every translation of each template can be chosen so
that its $f$-image is countably $m$-rectifiable and has the area formula
with tangential derivative, counted with the parametrization's
multiplicity.  The same conclusions hold on all lower strata needed
in a finite stratification.  At a generic intersection with a
complementary-dimensional affine test stratum, the source trace has a
regular manifold stratum and regular full-rank parameter branches,
with locally constant multiplicities.  Exceptional image values,
including branch junctions and dimension-drop intersections, may be
excluded from such intersections.
\end{lemma}

\begin{proof}
For $m=1$, Sobolev slicing gives absolute continuity on almost every
parallel line, and hence the one-dimensional area formula and Lusin
property.  A smooth curve chart can be foliated by transverse
translates and straightened by a local ambient smooth diffeomorphism,
giving the same conclusion in that chart.  For $m=2$ use
\cite[Theorem~1.3]{CsornyeiHenclMaly2010}: a $W^{1,2}$ homeomorphism
in dimension three has the two-dimensional Lusin property on almost
every parallel plane.  Together with Sobolev slicing and the usual
Lipschitz decomposition of a Sobolev map
\cite[Theorem~3(1)]{BojarskiHajlasz1993}, this gives the surface area
formula. Apply the Cs\"ornyei--Hencl--Mal\'y theorem after straightening a smooth graph
by an ambient smooth diffeomorphism.  Fixing any other translation
coordinates and then applying Fubini proves the assertion for the
translated graph templates.

Stratify the semialgebraic source parametrizations $P:K\to\Omega$,
their source images $P(K)$, and the ranks and overlaps of these
parametrizations, using a common finite refinement as in
\cite[Proposition~2.3 and Lemma~2.5]{HardtLambrechtsTurchinVolic2011}.
This stratification concerns $P$, not the generally non-semialgebraic
map $f\circ P$. On a full-rank branch the preceding graph argument applies,
and its parameter derivative is $Df(P)DP$.  On a lower-rank piece,
use the slicing assertion for its lower-dimensional image stratum;
for instance a one-dimensional image is rectifiable by the curve
assertion and has zero two-dimensional measure.  These observations
also make images of parameter-null sets negligible on the regular
branches.  The chain rule and tangential integrability may alternatively
be obtained by approximating $f$ by smooth maps in ambient Sobolev norm
and using Fubini to pass along a subsequence in parameter Sobolev norm.

We use the multiplicity area formula and the rectifiable-set coarea
formula in \cite[Chapter~2, Section~3, and Chapter~3, Section~2]{Simon2014Draft}.
The area formula and translation coarea imply that almost every
complementary affine test avoids the exceptional image sets.  At a
remaining source value, stratification, local inversion on regular
branches, and compactness leave a single local smooth source stratum
with a fixed finite number of parameter branches.  Dimension-drop
intersections of different pieces are negligible by the same argument.
More explicitly, full-rank strata of the finite refinement are
diffeomorphic onto their images, which can be made identical or
disjoint; their number, rather than compactness alone, bounds the
branch multiplicity. On relatively compact branch charts $P$ is
bi-Lipschitz, so parameter-null sets first have null source image and
then null target image by the separate sliced Lusin property.
Finally, $p\ge m$ in the admitted cases. Translation and Fubini give
finite tangential $m$-mass on almost every compact template, and the
area formula for its orthogonal projection makes the complementary
intersection count finite for almost every affine test.
The locally finite family involves only finitely many templates on
each compact set; intersect their full-measure choices and then use
a countable exhaustion.  This proves all the simultaneous assertions.
\end{proof}

We describe the target meshes used below.  Their shape need not be
uniform; only bounded local site counts and upper size bounds are
needed for target budgets.  Uniform source shapes will be obtained
separately in Lemma~\ref{lt:source-mesh}.

\begin{lemma}[Target mesh with lattice patches]
\label{lt:target-mesh}
In an open subset $V\subset\R^3$, prescribe finitely many separated
compact patches with buffers.  In a deep part of each patch one may
require a rotated rectangular lattice of side lengths $d$ or
$\eta d$, where $0<\eta\leq1$ is fixed.  There are locally finite
triangulations of $V$ agreeing there with the chain, or Kuhn,
triangulation of that lattice and with arbitrarily small prescribed
local upper size bounds.  For each local candidate, its sampling variables admit a physical
common translation $t$ and normalized relative coordinates $z$ whose
\emph{joint} density satisfies
\begin{equation}
 q(t,z)\leq Cd^{-3},\qquad
 t\in B(t_0,Cd),\quad z\in Z\subset B(0,C),
 \label{lt:mesh-joint-density}
\end{equation}
where the dimension and measure of $Z$ are bounded.  Deterministic
lattice relative coordinates are held fixed; when there are no free
relative coordinates their integration means a unit point mass.
The bounds are universal away from enlarged lattice patches and may
depend on $\eta$ inside them.  Only boundedly many candidates occur
within a bounded number of local mesh lengths.  Formula
\eqref{lt:mesh-joint-density} is used by joint integration, without
normalizing a conditional translation law after fixing $z$.
\end{lemma}

\begin{proof}
Choose a positive size function $d(u)$ with sufficiently small absolute
Lipschitz constant, satisfying the prescribed local upper bounds and
$d(u)\ll\dist(u,\partial V)$.  Make it constant, with the desired
common value, near each lattice patch.  Outside deep lattice regions
take a maximal packing at this scale, then independently jitter each
retained site in a ball of a fixed small fraction of its local scale,
with uniformly bounded density.  In a lattice patch use a common
uniform positional shift over one period.  Truncate the unshifted
lattice lists in their buffers, leaving overlap with the background
cloud; the separated lattice patches remain many steps apart.
The resulting sites cover at distance at most $Cd(u)$ and have
bounded counts within that distance.  A minimum separation between
every pair of sites is unnecessary.

Use their Delaunay subdivision, with a fixed pulling order to
triangulate nonsimplicial cells.  Its localization can be checked
inside the open domain.  An empty ball containing or passing through
a point near $u$ cannot have radius much larger than $d(u)$: moving
several mesh lengths from that point toward the center stays inside
$V$ and gives a location well inside the ball; the covering property
then puts a site strictly inside it.  The slow variation of $d$
justifies using the same scale during this argument.  Finite
restrictions of the cloud therefore stabilize near $u$ once they
contain a sufficiently enlarged local cloud.  Their ordinary Delaunay
subdivisions cover and agree there.  This proves local existence,
local finiteness, and compatibility of the subdivisions on $V$.
Every candidate uses a bounded list of neighboring sites at comparable
scales; coincidences occur with probability zero.

Deep in a rectangular lattice, the Voronoi description shows that the
Delaunay cells are the rectangular boxes.  Increasing lexicographic
pulling gives their Kuhn simplices.  Taking the scale small enough
places any prescribed deep compact set and all its touching simplices
within this region before sampling.

For the density assertion, first fix a candidate tuple of site
indices, before imposing its Delaunay selection event.  If it contains
lattice sites, they all belong to one shifted lattice.  Use its
physical shift as $t$ and write each unstructured site as
$x_0+t+dz_j$, with $x_0$ a fixed candidate anchor.  The original
joint law of the shift and the $k$ unstructured sites has density
at most $Cd^{-3(k+1)}$.  Replacing the latter positions by $z_j$
has Jacobian $d^{3k}$, giving \eqref{lt:mesh-joint-density}.
If no lattice site is present, use one site as $x_0+t$ and express
the other $k-1$ sites relative to it in units $d$; the Jacobian
$d^{3(k-1)}$ gives the same bound from the original density
$Cd^{-3k}$.  The normalized relative coordinates lie in a bounded
set because all candidate sites have comparable scales and lie in
one bounded-scale neighborhood.  The Delaunay selection event is
an indicator bounded by one and is discarded for an upper estimate.
All subsequent calculations integrate the unnormalized translation
slice for each $z$, then integrate $z$ over this bounded set.  Thus
no bound on a normalized conditional translation density is needed.
\end{proof}

Counting intersections with complementary cells is also central to
polyhedral deformation; see \cite[Proposition~2.2]{White1999}. Here
the weights count unsigned parameter visits, and the concentration
will be realized by homeomorphisms of the target.

\begin{definition}[Tie budget]
\label{lt:tie-budget}
In a tetrahedron $\sigma$ with vertices $v_0,\ldots,v_3$, write its
barycentric coordinates as $\lambda_i$.  For a set $I$ of $m+1$
vertices, the associated dual stratum is
\[
 Z_{\sigma,I}=
 \{\lambda_i=\max_j\lambda_j\ (i\in I),\quad
           \lambda_j<\max_i\lambda_i\ (j\notin I)\}
          \cap\relint\sigma.
\]
Write a source parametrization as $P:K\to\Omega$ and its measured
image as $\Gamma=f\circ P$; lengths and areas of $\Gamma$ mean
integrals on $K$.  For such a generic measured $m$-trace, its budget
$\mathcal B_m(\Gamma)$ is the sum, over its visits to these strata, of
\[
 \cH^m(\conv\{v_i:i\in I\}),
\]
with visits counted with parameter multiplicity.  A weighted budget
uses in addition the supremum of a prescribed nonnegative continuous
weight $w$ over $\sigma$.  A fixed linear change in the output may
also be applied when computing the flat weights.
\end{definition}

A triangle gives a simple model for this budget rule. Let
$v_0,v_1,v_2$ be its vertices and let $\lambda_i>0$ be interior
barycentric coordinates. For a constant $a>0$, consider the map
\[
 \sum_{i=0}^2\lambda_i v_i
 \longmapsto
 \frac{\sum_{i=0}^2\lambda_i e^{a\lambda_i}v_i}
      {\sum_{i=0}^2\lambda_i e^{a\lambda_i}}.
\]
As $a\to\infty$, a point with one largest coordinate moves toward
that coordinate's vertex. A regular curve in a sufficiently small
neighborhood of a point of $\lambda_0=\lambda_1>\lambda_2$, crossing
that tie once transversely with endpoints in the
two adjacent one-leader regions, instead runs between $v_0$ and
$v_1$. The concentration calculation below shows that one such
crossing contributes the edge length $|v_1-v_0|$ in the limit;
Figure~\ref{lt:tie-model} illustrates this local situation. In a
tetrahedron the same principle pairs a curve with a two-leader
surface and an edge weight, or a surface with a three-leader line
and a triangle weight. This is why the budget uses complementary
dimensions and flat primal faces.

\begin{figure}[tb]
\centering
\begin{tikzpicture}[x=0.95cm,y=0.95cm,font=\small]
 \begin{scope}
  \fill[blue!5] (0,0)--(2,0)--(2,1.155)--(1,1.732)--cycle;
  \fill[orange!9] (4,0)--(3,1.732)--(2,1.155)--(2,0)--cycle;
  \draw (0,0)--(4,0)--(2,3.464)--cycle;
  \draw[dashed,gray] (2,0)--(2,1.155)--(1,1.732);
  \draw[dashed,gray] (2,1.155)--(3,1.732);
  \draw[blue!65!black,thick,-{Stealth[length=2mm]}]
    (1.2,0.57) .. controls (1.7,0.72) and (2.25,0.45) .. (2.8,0.62);
  \node[left,blue!65!black] at (1.2,0.57) {$\Gamma$};
  \node[below left] at (0,0) {$v_0$};
  \node[below right] at (4,0) {$v_1$};
  \node[above] at (2,3.464) {$v_2$};
  \node[font=\scriptsize,fill=white,inner sep=1pt] at (2,0.95)
    {$\lambda_0=\lambda_1>\lambda_2$};
 \end{scope}
 \draw[-{Stealth[length=2mm]}] (4.55,1.6)--(6.25,1.6)
    node[midway,above] {$a\to\infty$};
 \begin{scope}[xshift=7.1cm]
  \draw[gray!60] (0,0)--(2,3.464)--(4,0);
  \draw[blue!65!black,very thick,-{Stealth[length=2mm]}] (0,0)--(4,0);
  \node[below left] at (0,0) {$v_0$};
  \node[below right] at (4,0) {$v_1$};
  \node[above] at (2,3.464) {$v_2$};
  \node[below] at (2,-0.3) {limiting edge contribution};
 \end{scope}
\end{tikzpicture}
\caption{A two-dimensional model of a tie budget. Dashed segments
are the two-leader ties. The curve crosses one of them away from
the three-leader center; its image concentrates on the corresponding
edge. The right panel records the limiting contribution, not the
image under a finite-parameter homeomorphism.}
\label{lt:tie-model}
\end{figure}

For generic samples all relevant hits are finite on compact measured
pieces.  Higher ties, simplex-boundary junctions, and lower-dimensional
source exceptions are avoided by Lemma~\ref{lt:trace-slicing} and
translation coarea.  Counts consequently concern isolated hits at
regular underlying source strata; no differentiability of the original
image surface at such a hit is being assumed.

\subsection{A variable simplex squeeze}

\begin{lemma}[Compatible simplex squeezing]
\label{lt:squeeze}
Let $\mathcal T$ be a locally finite triangulation of a three-dimensional
open domain.  Assign constants $0<b_\sigma\leq1$ to its tetrahedra,
and to every nonempty face $\rho$ assign the minimum of the constants
of its incident tetrahedra.  On each tetrahedron put
\begin{align}
 b(\lambda)&=
 \min_{\varnothing\ne\rho\subseteq\sigma}
 \left(b_\rho+
        \frac{\max_{i\notin\rho}\lambda_i}{\max_j\lambda_j}\right),
       \label{lt:b-definition}\\
 a(\lambda)&=e^{1/b(\lambda)}-e,
 \qquad
 T_i(\lambda)=
 \frac{\lambda_i e^{a(\lambda)\lambda_i}}
             {\sum_j\lambda_j e^{a(\lambda)\lambda_j}},
       \label{lt:squeeze-definition}
\end{align}
where the maximum over an empty set in \eqref{lt:b-definition} is zero.
These formulas define a face-preserving, locally bi-Lipschitz
self-homeomorphism $T$ of the domain, with finitely many $C^1$ formulas
on closed pieces of every compact subset.  The normalized Lipschitz
constant of $b$ is bounded absolutely, independently of its minimum,
and a.e.
\begin{equation}
 |Da|\leq C(a+1)(1+\log(a+1))^2.
 \label{lt:a-derivative}
\end{equation}
If all the incident constants in a mesh neighborhood are one, $T$ is
the identity on its inner simplices.

Let $G$ parametrize finitely many compact $m$-traces, $m\in\{1,2\}$,
by finite closed-piece $C^1$ formulas, and suppose their intersections
with the $Z_{\sigma,I}$ are regular transverse hits away from higher
ties and simplex boundaries.  Every counted hit is in the interior
of each counted regular $m$-dimensional parameter branch; parameter
boundaries and lower-dimensional parameter strata avoid the counted
dual strata.  As the lower bounds of $a$ tend to
infinity on the measured neighborhood,
\begin{equation}
 \int J_mD(T\circ G)
 \longrightarrow
 \sum_{\text{hits on }Z_{\sigma,I}}
       \cH^m(\conv\{v_i:i\in I\}),
 \label{lt:squeeze-area}
\end{equation}
with multiplicity.  The local convergence, and in particular its upper
bound with arbitrarily small error, is uniform when other floors of
$a$ are made arbitrarily larger.  Continuous weights may be bounded
by their simplex suprema in this conclusion.
\end{lemma}

\begin{proof}
Since the candidate $\rho=\sigma$ gives $b\leq b_\sigma\leq1$,
$a\geq0$.  Every denominator $\max\lambda_j$ is at least $1/4$,
so the normalized Lipschitz constant of $b$ is absolute.  The chain
rule gives \eqref{lt:a-derivative}.  On a face $\tau$, replacing a
candidate $\rho$ by $\rho\cap\tau$ leaves its nonzero outside
coordinates unchanged and can only lower its constant, because
$b_{\rho\cap\tau}\leq b_\rho$.  An empty intersection gives an
outside-coordinate ratio of one and cannot improve on $b_\tau\leq1$.
Thus restrictions agree on faces.  Positivity, zero coordinates, and
the order of coordinates are preserved by \eqref{lt:squeeze-definition}.

We verify invertibility even though $a$ is variable.  Given output
ratios $z_i=T_i/T_{\max}\in[0,1]$, for a trial $D\geq0$ let $r_i$
be the unique solution of
\begin{equation}
 z_i=r_i e^{-D(1-r_i)},\qquad 0\leq r_i\leq1.
 \label{lt:ratio-inverse}
\end{equation}
The function of $r_i$ on the right is strictly increasing.  Each
$r_i$ is nondecreasing in $D$, with the zero and maximal ratios fixed
at zero and one.  Put $R=\sum_i r_i$ and $\lambda_i=r_i/R$.
The trial value of $a$ is $DR$, an increasing function of $D$ with
increase at least that of $D$.  In \eqref{lt:b-definition} the
outside-coordinate ratios are exactly $r_i$; hence $b$ at this trial
inverse is nondecreasing in $D$.  Its prescribed value
$e^{1/b}-e$ is nonincreasing.  The equation
\[
 DR=e^{1/b(\lambda)}-e
\]
has a unique solution, by continuity, its sign at $D=0$, and its
positive sign for large $D$.  This proves bijectivity.
On an individual closed simplex $b$ has a positive lower bound, so
the solution has bounded $D$.  Equation~\eqref{lt:ratio-inverse}
has uniformly Lipschitz inverse in this bounded range, including at
zero output coordinates.  The scalar equation's left-minus-right
increase is at least the increase in $D$, so its root depends
Lipschitz-continuously on the output.  Since $T_{\max}\geq1/4$,
forming the output ratios is Lipschitz as well.  The forward formula
is Lipschitz on that simplex, and its inverse is Lipschitz.  The
finite min/max alternatives in \eqref{lt:b-definition} give finitely
many closed pieces with $C^1$ extensions of the formulas.  Local
finiteness proves all the global assertions.  If the incident
constants are one, the minimum is one, so $a=0$ and $T=\Id$.

It remains to prove the concentration assertion.  On a compact trace
set away from $(m+1)$-fold leader ties, at most $m$ coordinates can be
near the maximum.  Coordinates outside that set and all their
derivatives are exponentially small in $a$, multiplied only by
powers of $a$ and $\log(a+e)$ by \eqref{lt:a-derivative}.  The active
coordinates have total mass one up to those errors and therefore
take values in a face of dimension at most $m-1$.  Their $m$-Jacobian
vanishes.  Expanding the actual $m$-Jacobian then shows uniform
convergence to zero away from the counted ties, including along
simplex faces by the same piecewise calculation.

Near a transverse hit relabel the leaders $0,\ldots,m$ and use
\[
 u_i=\lambda_i-\lambda_0,\qquad 1\leq i\leq m,
\]
as trace coordinates.  All leader coordinates have a fixed positive
lower bound there, and the remaining coordinates are separated below
them.  Normalize the probabilities on the leaders first.  Their
logit differences are
\[
 a(u)u_i+\log\frac{\lambda_i(u)}{\lambda_0(u)}.
\]
Inactive probabilities and derivatives are exponentially small.
Write $a_0=a(0)$.  Since
$b(0)=1/\log(a_0+e)$ and $b$ is uniformly Lipschitz, on
$u=t/a_0$ with bounded $t$ one has
\[
 \log\frac{a(u)+e}{a_0+e}
       =O\left(\frac{(\log(a_0+e))^2|t|}{a_0}\right)=o(1).
\]
The scaled derivative terms coming from $Da$ tend to zero by
\eqref{lt:a-derivative}.  The leader map and its a.e. derivatives
therefore converge locally uniformly, apart from its harmless
piece seams, to
\[
 t\longmapsto
 \left(\frac{e^{t_i}}{1+\sum_{j=1}^m e^{t_j}}\right)_{i=1}^m.
\]
This softmax map is a diffeomorphism from $\R^m$ onto the interior
of the standard $m$-simplex.  Multiplication by the physical edge
vectors consequently gives exactly the flat $m$-area of the leader
face, provided the tails vanish.

Here are uniform tail bounds.  The determinant of the active
probability derivative is its probability product times the
determinant of the logit derivative.  The logit derivative has the
form $aI+u\otimes Da+B$, with $B$ bounded for the fixed trace chart.
The probability product is at most $C e^{-ca|u|}$.  Thus the
absolute determinant, modulo exponentially small inactive terms,
is bounded by
\begin{equation}
 C(1+a+|Da|\,|u|)^m e^{-ca|u|}.
 \label{lt:tail-bound}
\end{equation}
For $|u|\leq a_0^{-1/2}$, the preceding logarithmic estimate is
uniform and gives $a/a_0\to1$ there.  After scaling, the
determinant has an integrable bound $Ce^{-c|t|}$, because
$|Da||u|/a_0\leq C(\log a_0)^2a_0^{-1/2}=o(1)$.
For
\[
 a_0^{-1/2}<|u|\leq c_1/\log(a_0+e),
\]
choose $c_1$ small relative to the fixed Lipschitz bound for $b$.
Then $a\geq a_0^{3/4}$ for large $a_0$, and
$a|u|\geq a_0^{1/4}$.  Exponential decay in
\eqref{lt:tail-bound} absorbs all polynomial and logarithmic
factors uniformly for every larger $a$.
In the remaining sufficiently small neighborhood,
\[
 b(u)\leq b(0)+C|u|\leq C'|u|,
 \qquad a(u)+e\geq e^{c'/|u|}.
\]
For small $r=|u|$, the expression $a^N e^{-car}$ is decreasing
once $a\geq e^{c'/r}/2$, for any fixed $N$.  Increasing $N$ to
absorb the logarithms in \eqref{lt:tail-bound} therefore gives an
integrable bound of the form
\[
 C\exp(Nc'/r)\exp(-c''r e^{c'/r}),
\]
independent of how high the local floors are.  At every fixed
$u\ne0$ the integrand tends to zero as the floor tends to infinity.
Dominated convergence treats this last region.  These three bounds
prove tail vanishing and \eqref{lt:squeeze-area}.

For a compact measured list there are only finitely many charts and
hits.  One may therefore prescribe the floors simultaneously, with
any allocated small errors.  Local finiteness then permits such
choices for a locally finite list, including summable errors.  The
map moves a point only inside its own simplex, so multiplication by
simplex suprema of continuous weights proves the weighted upper
bound.  The same reasoning applies to every regular parameter
branch, and hence respects multiplicity.
\end{proof}

\subsection{Transfer and selection of the budgets}

\begin{proposition}[Topological transfer of trace budgets]
\label{lt:budget-transfer}
Fix generic trace data as in Lemma~\ref{lt:trace-slicing} and a
sufficiently fine generic mesh from Lemma~\ref{lt:target-mesh}.
There is an onto locally bi-Lipschitz homeomorphism
$H_0:\Omega\to\Omega'$, locally specified by finitely many
$C^1$ formulas on closed pieces, whose image lengths and areas on
the measured traces are at most their tie budgets plus arbitrarily
small prescribed local errors.  Nonnegative continuous output
weights are allowed, with the simplex-supremum budgets of
Definition~\ref{lt:tie-budget}.

The map can be arbitrarily close to $f$, with reciprocal closeness
on interior compact sets, within the mesh displacement and chosen
local modification errors.  It can retain smaller exact affine
germs previously installed at a locally finite family of crossings
and finitely many star centers, and can retain strict compact
forward and inverse safety conditions, including angular cone
conditions, whenever they have positive margins.  For this relative
usage, measured data are separated from buffered open exemptions
around the protected germs: a measured trace either avoids the
exempt cores, or its weight vanishes there with a buffer.  Every
germ required to remain exact in $H_0$ is contained in one of these
buffered exemptions.  The affine crossing germs installed at the
newly cleaned measured hits are retained only in the intermediate
map $G$ for the concentration calculation; no affine-germ condition
at those hits is imposed on $T\circ G$.  The mesh and all measured
configurations are fixed before the squeezing floors are selected.
\end{proposition}

\begin{proof}
Each area-line hit is an isolated intersection of the image of a
regular source surface with an affine target interval.  Apply
Lemma~\ref{lt:clean-crossing}, then Corollary~\ref{lt:affine-crossing},
in a paired source and target neighborhood.  For a curve-plane hit,
interchange source and target and clean the preimage plane.  The
supports can be chosen disjoint and locally finite.  Area-line hits
avoid curve-plane hits and the measured curve images; curve-plane
changes can avoid all area-test lines.  If parametrizations share a
regular source stratum at a hit, the operation is performed once on
that stratum and its fixed multiplicity is retained.  Lower strata
and other sheets are absent there by genericity.  A surviving hit
remains in its intended source stratum and target tie simplex, and
the cleaning does not increase its count.

Use Lemma~\ref{lt:relative-germs} to obtain a locally bi-Lipschitz
map $G$ with finite closed-piece $C^1$ formulas, keeping smaller
exact crossing and center neighborhoods.  Take the approximation
smaller than the positive gaps from all forbidden strata off those
neighborhoods.  The measured intersections of $G$ are now regular
transverse intersections and their counts do not exceed the
original budget counts.  Choosing all supports and approximation
errors finely also retains the stated compact forward and inverse
conditions.  These choices use only topology and qualitative
approximation; no derivative bound for $G$ is charged to a budget.

Apply Lemma~\ref{lt:squeeze} in the target and set $H_0=T\circ G$.
It has the asserted local regularity and is an onto homeomorphism.
On each compact measured list choose the floors high enough that
the image measurements cost at most the flat weights plus their
allocated errors.  Finitely many such lists meet each simplex, so
the uniformity in arbitrarily higher floors permits simultaneous
choices for the locally finite family.

For the relative assertion, first keep each protected-germ modification
deep inside its exempt region.  Make the mesh so fine that all
simplices touching the protected no-movement core, and their needed
neighbors, lie inside that exemption.  Assign the value one to
their $b_\sigma$ constants; then $T=\Id$ on a smaller neighborhood
of the germ.  All simplices required for measured points, including
points where the final weight is nonzero, remain buffered away
from these cores and use the original $f$-crossings.  Subsequent
cleaning supports are disjoint from the core neighborhoods.
Every motion by $T$ stays inside a simplex.  Thus the chosen fine
mesh retains the strict margins and the reciprocal compact
accuracies, and $w(TG(x))\leq\sup_\sigma w$ whenever $G(x)\in\sigma$.
This proves the weighted estimate as well.  No unweighted estimate
through an exempt affine-star core is asserted or needed.
\end{proof}

\begin{lemma}[Expected budgets]
\label{lt:budget-expectation}
For an $m$-trace with finite image measure, the meshes in
Lemma~\ref{lt:target-mesh} satisfy
\begin{equation}
 \mathbb E\,\mathcal B_m(\Gamma)
             \leq C\int_K J_mD(f\circ P),
 \label{lt:expected-budget}
\end{equation}
with parameter multiplicity understood.  The constant is universal
outside enlarged lattice patches and may depend on their fixed
aspect ratios inside them.  The estimate holds locally with
comparable enlargements.  For continuous weights it holds with
the weighted trace integral and an arbitrarily small enlargement
error as the local mesh sizes tend to zero.  Fixed linear output
normalizations may be included in all areas.
\end{lemma}

\begin{proof}
Use the joint variables $(t,z)$ of
\eqref{lt:mesh-joint-density} for a candidate tetrahedron.
For each fixed $z$, its dual $(3-m)$-piece $Z_z$ has diameter
$O(d)$ and measure at most $Cd^{3-m}$; translation by $t$
produces the actual dual piece.  If $N(t,z)$ is its intersection
count with the measured trace in the candidate's fixed enlarged
neighborhood, translation coarea gives
\[
 \int N(t,z)\,dt
 \leq C d^{3-m}\,
          \text{local trace image $m$-measure}.
\]
Indeed this is the area formula for the difference of the trace
and dual-piece parametrizations, whose angle factor is at most one.
The joint density bound therefore gives, for every such $z$,
\[
 \int q(t,z)N(t,z)\,dt
 \leq C d^{-m}\,
          \text{local trace image $m$-measure}.
\]
These are unnormalized translation integrals.  Integrating $z$
over its bounded normalized set only changes the constant.
The flat $m$-face weight is at most $Cd^m$ for every $z$, so it
cancels the remaining scale.  Discard the Delaunay selection
indicator for this upper estimate.
Only boundedly many candidate tuples occur locally and their
enlargements have bounded overlap, proving
\eqref{lt:expected-budget}.  The same coarea calculation gives
generic avoidance of negligible sets and finiteness of compact
counts.  A continuous weight's supremum on each vanishing simplex
differs from its value at a measured point by its local modulus
of continuity.  On a compact set multiply that modulus by the
finite trace measure; a locally finite exhaustion and allocated
errors give the weighted assertion.
\end{proof}

\begin{lemma}[Length budgets with arbitrarily high probability]
\label{lt:curve-probability}
For any compact rectifiable measured curve $\Gamma$ of finite
length, any $\delta>0$, and any $e>0$, sufficiently small local
mesh upper bounds make
\begin{equation}
 \mathcal B_1(\Gamma)\leq C\operatorname{length}(\Gamma)+e
 \label{lt:curve-high-probability}
\end{equation}
hold with probability at least $1-\delta$.  The constant $C$ is
independent of $\delta,e$ and the required mesh resolution; it is
universal away from enlarged patches.  Prescribed summable failure
probabilities and local errors can be imposed simultaneously on a
countable locally finite family of curves.  In a deep anisotropic
lattice patch there is also a grid-unit version: rescale to unit
grid aspect, estimate length there, and convert the weight back
using the longest physical side, with a constant independent of
the aspect ratio in those units.
\end{lemma}

\begin{proof}
Discard an arbitrarily small amount of length with multiplicity
and decompose the remainder into finitely many compact pieces on
$C^1$ graphs of arbitrarily small slope $s$ over straight axes.
Such pieces may be taken injectively from the regular parameter
branches, so multiplicity is accounted for by separate pieces.
Cover their axis projections by finitely many intervals of total
length at most the sum of their graph-piece lengths plus an
arbitrarily small error.  Extend the graphs over these intervals
with comparably small slopes.

For each fixed relative-coordinate value $z$ in the joint
integration formula \eqref{lt:mesh-joint-density}, classify a dual
plane by its unit normal $\nu(z)$ and the graph axis $e$.  If
$|\nu\cdot e|>2s$, the restriction of its affine defining function
to the graph is strictly monotone.  Each such plane patch meets
the graph over its interval at most once.  Charge a contributing
tetrahedron to an axis interval of length comparable to $d$
centered at a hit.  Intersecting charged intervals have comparable
scales, by slow grading, and their vertices belong to a bounded
local candidate list.  Thus the charges have bounded overlap.
The sum of mesh diameters, and hence of the flat edge weights,
is bounded by a constant times the graph length plus the local
mesh upper bound at each interval endpoint.  After the finite
cover is fixed, all endpoint overhangs can be made as small as
desired by the mesh choice.

For the remaining normals, the translation-coarea factor on this
graph is at most $Cs$.  For each fixed $z$, integration over the
physical translation variable bounds the crossing integral by
$Cs d^2$ times the local graph length.  Multiply by the joint
density bound $Cd^{-3}$ and the edge weight $Cd$, and then
integrate over the bounded relative-coordinate set.  This gives
expected cost at most $Cs$ times the graph length.  The normal
classification is translation invariant, so it is legitimate
inside this unnormalized joint integral.  The omitted trace length has correspondingly
small expected cost by Lemma~\ref{lt:budget-expectation}.  Choose the omitted length
and $s$ sufficiently small, then choose the mesh bounds for the
finite graph cover.  Markov's inequality applies only to these
small remainders, and makes their failure probability below
$\delta$ without changing the constant on the retained graph
length.  This proves \eqref{lt:curve-high-probability}.

For a locally finite family assign summable failure probabilities
and errors and impose the corresponding local upper bounds
together.  A union bound proves the simultaneous assertion.
In a deep lattice patch perform the entire argument after the
anisotropic rescaling.  There are finitely many rational tie-plane
families in these units; the same argument, or translated Riemann
sums for their projected length coarea, gives the claimed constant.
Physical flat weights are bounded by the longest-side factor
times their grid-unit weights.
\end{proof}

\begin{lemma}[Sharp costs along a thin-grid subspace]
\label{lt:sharp-grid}
In a deep lattice patch adapted to a fixed $k$-plane $E$, use side
$d$ along $E$ and side $\eta d$ along $E^\perp$, where
$k\in\{1,2\}$.  For $m=k$, the translation-averaged tie-budget
coefficient on a simple $k$-vector tangent to $E$ is at most
$1+C\eta$ times its Euclidean norm.  For general measured
$m$-vectors, tangent contributions have constants uniform in
$\eta$, and errors have a finite constant $C_\eta$.
After a fixed linear output normalization $A$ which keeps the two
subspaces orthogonal, the sharp coefficient relative to the normalized
tangent length or area is at most $1+C_A\eta$. Here $C_A$ depends only on
$\|A\|_{\op}$ and $\|A^{-1}\|_{\op}$. The tangent bounds remain
uniform in $\eta$, while the error constants may also depend on these
two norms.
The grid-unit bound of Lemma~\ref{lt:curve-probability} also has no
aspect-dependent factor on tangent lengths.
\end{lemma}

\begin{proof}
Rescale the grid to unit cubes, temporarily recording physical
flat weights separately.  For $k=2$ slice a cube at a generic
height in its thin direction.  Each Kuhn tetrahedron projects to
a triangle with one repeated vertex.  Write the two unsplit
projected barycentric masses as $q_s,q_t$.  On a fixed-height
slice, either part $\lambda_0$ of the split mass is a constant
minus a sum of some of $q_s,q_t$, with coefficients zero or one.
Consequently
\[
 \det D_{(q_s,q_t)}(q_s-\lambda_0,q_t-\lambda_0)>0.
\]
These are also the orientation coordinates for a leader triangle
with three different projected vertices.  Thus every such
projected triangle is counted positively.

The total projected weight on a horizontal period is exactly
its area.  One way to see this without estimating individual
positions is to take a constant large $a$ in
\eqref{lt:squeeze-definition}.  Projection of this periodic
simplexwise map on the horizontal slice is a torus map with
periodic displacement and degree one.  Its signed Jacobian
integrates to one base area.  Lemma~\ref{lt:squeeze} converts
this integral, in the large-$a$ limit, into the signed projected
triangle weights.  Generic heights give only finitely many
transverse hits away from simplex boundaries; the signs just
computed are positive, and repeated projections contribute zero.
For a direct check, the two Kuhn tetrahedra whose vertical step
comes first supply total projected area one at heights
$2/3<z<1$, those whose vertical step is middle supply one for
$1/3<z<2/3$, and those whose vertical step is last supply one for
$0<z<1/3$.  Each nondegenerate leader triangle contributes $1/2$
at its applicable height.  The finitely many exceptional heights
do not affect translation averaging.

In physical units a nondegenerate projected triangle differs
from its projected area by at most $C\eta$ times the base area;
a triangle with repeated projections has area at most
$C\eta d^2$.  The number of hits per period is uniformly bounded.
The total physical area cost is therefore at most
$(1+C\eta)d^2$.

For $k=1$, fix generic coordinates transverse to the long axis
and traverse one axial period.  Within a Kuhn tetrahedron the
varying barycentric pair transfers mass from its lower axial
position to its upper axial position, while the other masses
are fixed.  A switch between leaders with different axial
projections is consequently always forward.  Two fixed
unsplit masses do not tie at a generic transverse coordinate.
The projected periodic map has degree one, so total projected
length is exactly the axial period.  Edges with equal axial
projection, and the transverse parts of the other edges, cost
only $O(\eta)$ in period units.  This proves the sharp length
coefficient.

Translation coarea expresses each coefficient on a trace
Jacobian as a finite homogeneous sum of absolute linear forms
on its $m$-vector.  The parallel-vector computations therefore
give the asserted tangent coefficients.  Splitting a vector
into its tangent part and its error and using the triangle
inequality gives a finite $C_\eta$ on the error.  A fixed
normalization preserving orthogonality of the two subspaces
changes these constants only through its two norm bounds.
The final length assertion follows by making the anisotropic
rescaling before applying Lemma~\ref{lt:curve-probability}.
\end{proof}

\begin{remark}[Order of choices for weighted budgets]
\label{lt:weighted-order}
Additional compact configurations and continuous weights may be
fixed before selecting the target mesh bounds.  They change how
fine the mesh must be for the preceding estimates, but not the
constants in those estimates.  Small open exemptions, their
buffered no-movement cores, and strict angular safety conditions
are included in this convention.  Squeezing floors are selected
only after the regular transverse map $G$ and its compact
measurement lists are fixed.
\end{remark}

\subsection{Source meshes with uniform shapes}

\begin{lemma}[Shape-controlled source sampling]
\label{lt:source-mesh}
There are locally finite source meshes of uniform shape with
arbitrarily small local upper size bounds, preserving finitely
many separated deep cubical lattice patches in prescribed
orientations and with prescribed common fine scales.  Deep
cells may remain cubes; elsewhere they are tetrahedra, with
matching facet diagonals at interfaces.  For an integrable
nonnegative function $g$ and $m\in\{1,2\}$, their sampled
$m$-faces satisfy the local estimate
\begin{equation}
 \mathbb E\sum_F \ell_F^{3-m}\int_F g\,d\cH^m
                    \leq C\int g,
 \label{lt:source-fubini}
\end{equation}
with comparable local enlargements and bounded overlap understood.
In particular this applies to $g=|Df|^p$ and controls the
tangential derivative trace integrals.  Lattice shifts can also
be uniform over any prescribed macro-period formed by grouping
fine cubes.
\end{lemma}

\begin{proof}
Use the construction of Lemma~\ref{lt:target-mesh}, now with
cubical rather than anisotropic deep patches, and impose
uniform altitude avoidance on the unstructured sites.  Conditional
on the lattice shifts, enumerate these sites and place them
successively.  When placing a site, exclude a neighborhood of
thickness $\varepsilon_0$ times its scale around the affine span
of every nearby tuple of at most three earlier or lattice
vertices.  Candidate neighbor lists are fixed and bounded by
localization.  A neighborhood of each such plane, line, or
point removes at most $C\varepsilon_0$ of the jitter volume.
For an absolute sufficiently small $\varepsilon_0$, at least
half the jitter volume remains.  Choose uniformly in the
remaining portion.

Any affinely independent lattice prefix of a candidate simplex
comes from one rotated cubical lattice.  There are only finitely
many relevant normalized configurations, so its nonzero
determinant or lower-dimensional volume has an absolute lower
bound.  Each subsequently placed vertex has altitude at least
$\varepsilon_0$ times scale over the affine span of the preceding
vertices.  Induction gives a uniform lower bound on every
nondegenerate simplex determinant after scaling.  Together with
the upper diameter bound this gives uniform shapes.  Fully
lattice simplices have the usual finite collection of Kuhn
shapes.  Deep cells may be left cubical, using the same
diagonals on shared facets when meeting triangulated cells.
Block modifications can consequently be kept inside cubical
regions, with whole shared squares on their outer boundaries.

Given the previously placed sites and the lattice shifts, each
new jitter has conditional density at most twice its original
uniform density.  Successive integration of the unselected sites,
or the corresponding iterated-expectation bound on rectangular
events, therefore bounds the joint density of any fixed tuple of
$k$ selected unstructured sites by $C^k d^{-3k}$, uniformly in
the lattice shifts.  The jitters need not be independent.  Include the uniform density of a candidate's lattice shift, if
present, and make the same change of variables as in
Lemma~\ref{lt:target-mesh}.  This again gives a joint density
$q(t,z)\leq Cd^{-3}$ for physical translation and normalized
relative coordinates.  Discard any later selection indicator.
For each fixed $z$, an $m$-face has area $O(d^m)$, and Fubini gives
\[
 \int dt\int_{t+F_z}g\,d\cH^m
       \leq Cd^m\int_{U}g,
\]
where $U$ is the candidate's comparable enlarged neighborhood;
only translations in its bounded support are included on the
left.  Multiplying by $Cd^{-3}$, then integrating $z$ over its
bounded set, gives $Cd^{m-3}\int_U g$ for the expected face
integral.  Its mesh scale is comparable to $d$ by shape control.
Multiplication by $\ell_F^{3-m}$, followed by bounded candidate
counts and overlap, proves \eqref{lt:source-fubini}.

For grouping fine cubes into macro-cubes, choose independently
a uniform fine-period shift and a uniform discrete grouping
offset.  Their sum is uniform over the macro-period.  Truncation
of a lattice list concerns sites inside its fixed buffered patch;
it does not translate the transition-site lists by many mesh
steps.  Thus every mixed local candidate still uses only one
fine-scale physical translation variable after subtracting that
shift and normalizing its jitter coordinates.  Its bounded joint
density, and hence the unnormalized joint Fubini calculation,
remain unchanged.
\end{proof}
